\documentclass{article}
\usepackage[a4paper, total={6in, 8in}]{geometry}
\usepackage{graphicx} % Required for inserting images
\usepackage{amsmath,amssymb,xcolor,bm,amsthm,mathtools,comment, enumerate}
\newcommand{\eps}{\varepsilon}
\usepackage{tikz}
\newtheorem{theorem}{Theorem}[section]
\newtheorem{proposition}[theorem]{Proposition}
\newtheorem{remark}[theorem]{Remark}
\newtheorem{definition}[theorem]{Definition}
\newtheorem{lemma}[theorem]{Lemma}
\newtheorem{corollary}[theorem]{Corollary}
\newtheorem{example}[theorem]{Example}
\newtheorem{assumption}[theorem]{Assumption}
\newtheorem{conjecture}[theorem]{Conjecture}
\newtheorem{question}{Question}

\title{Doubling an admissible hypercube matching profile into a perfect admissible tuple}
\author{Joshua Erde}
\date{arxiv:2404.03950v3; CC BY 4.0}
\begin{document}
\maketitle
\noindent\textit{Editorial note.} The counted unit is the article's own Lemma l:admissible at source lines 67-69 together with the proof environment that begins on the next line (70-84); the pairing is the ordinary adjacent one, and the proof ends before the first figure of the article at source line 86.\par\medskip
\noindent\textit{Editorial note.} The retained context at source lines 30-36 and 51 is the article's own statement of the Balister-Gyori-Schelp necessary condition and conjecture, and the article's own definitions of profile, admissible, even and perfect tuples. The retained passage starts at the identity at source line 30 rather than at the preceding sentence at source line 29, which carries a bibliographic citation that is not part of the retained fragment; no line of the retained text was altered.\par\medskip
\noindent\textit{Editorial note.} The native preamble is retained from source lines 2-18, that is the article class, geometry, AMS, bm, mathtools, comment, enumerate and tikz packages and the author's theorem declarations; the title and author lines at source lines 19-20 are not part of the preamble fragment.\par\medskip
\noindent\textit{Editorial note.} No retained line contains a citation, so the material carries no bibliography; the retained lines refer only to the label c:BGS, which the retained conjecture defines.\par\medskip
\noindent\textit{Editorial note.} The article's Theorem t:perfect, its other lemma l:first and the discussion theorem c:admissible are not counted here and are not used as premises.\par\medskip
\noindent\textit{Editorial scope.} One original result: the article's Lemma l:admissible with its complete author proof, plus the article's own definitions of matching profiles, admissible, even and perfect tuples and the Balister-Gyori-Schelp necessary condition. The article's Theorem t:perfect and its other statements are not counted here and are not used as premises.\par\medskip
\noindent\textit{Allowed standard prerequisites (editorial).}\begin{itemize}
\item Profiles of matchings in the hypercube. For a matching M in Q\textasciicircum{}n the profile is the vector x in N\textasciicircum{}n whose i-th coordinate counts the edges of M whose endpoints differ in the i-th coordinate; a tuple is admissible when it is the profile of some matching, even when every coordinate is even, and perfect when the sum of the coordinates equals 2\textasciicircum{}\{n-1\}; the article's own definitions are retained at source line 51.. Fixes all objects of the counted lemma: admissible x, the perfect property of (2x, x\_\{n+1\}) and the count 2||x||\_1 + |X| = 2\textasciicircum{}n.
\item Balister-Gyori-Schelp partitioning problem. Given d\_1,...,d\_\{2\textasciicircum{}\{n-1\}\} in the vector space over F\_2 of dimension n, the sum of all d\_i must vanish if the cube is to be partitioned into pairs with those differences; the article records this identity and the conjecture that it is also sufficient at source lines 30-36, and the case of unit-vector differences is the matching-profile problem.. Background for why evenness is the necessary condition and why a perfect even profile is the interesting case; the counted lemma itself is used to build the inductive argument for that conjecture in the unit-vector case.
\end{itemize}\medskip
\begin{equation}\label{e:null}
\sum_{i=1}^{2^{n-1}} \bm{d_i} = \sum_{i=1}^{2^{n-1}}  \bm{a_i} - \bm{b_i} =   \sum_{i=1}^{2^{n-1}}  \bm{a_i} + \bm{b_i} = \sum_{\bm{v} \in \mathbb{F}_2^n} \bm{v} = \bm{0}.
\end{equation}
They conjectured that this necessary condition was also sufficient.
\begin{conjecture}[Balister, Gy{\H{o}}ri, and Schelp]\label{c:BGS}
Let $n \geq 2$ and let $\bm{d_1},\ldots, \bm{d_{2^{n-1}}}$ be elements of $\mathbb{F}_2^n$ such that $\sum_{i=1}^{2^{n-1}} \bm{d_i} = \bm{0}$. Then there exists a partition of $\mathbb{F}_2^n$ into pairs $\{\bm{a_i},\bm{b_i}\}$ such that $\bm{a_i} - \bm{b_i} = \bm{d_i}$ for each $i \in \left[2^{n-1}\right]$.
\end{conjecture}

Given a matching $M \subseteq Q^n$ we say the \emph{profile} of $M$ is the vector $\bm{x}=(x_1,\ldots, x_n) \in \mathbb{N}^n$ such that $M$ contains $x_i$ edges whose endpoints differ in the $i$th coordinate, and we are interested in which tuples $\bm{x} \in \mathbb{N}^n$ are achievable as the profile of some matching. We call such tuples \emph{admissible}. In this case, the obvious necessary condition for a tuple to be the profile of a perfect matching is that each unit vector appears an even number of times as a matching edge, i.e., each $x_i$ is even, in which case we say that $\bm{x}$ is \emph{even}, and Conjecture \ref{c:BGS} would imply that this condition is also sufficient. The aim of this short note is to verify the conjecture in this case. Let us call a tuple $\bm{x} \in \mathbb{N}^n$  \emph{perfect} if $||\bm{x}||_1 := \sum_i x_i = 2^{n-1}$.

\begin{lemma}\label{l:admissible}
If $\bm{x} \in \mathbb{N}^{n}$ is admissible and $x_{n+1} = 2^{n} - 2||\bm{x}||_1$, then $(2\bm{x},x_{n+1})$ is a perfect admissible tuple.
\end{lemma}

\begin{proof}
By assumption, there is a matching $M_1 \subseteq Q^{n}$ with profile $\bm{x}$. We form a matching $M_2 \subseteq Q^{n+1}$ by taking disjoint two copies of $M_1$, that is,
\[
M_2 := \{\big((\bm{u},0),(\bm{v},0)\big), \big((\bm{u},1),(\bm{v},1)\big) \colon (\bm{u},\bm{v}) \in M_1 \}
\]
and note that $M_2$ has profile $
(2\bm{x},0)$.

Let $X \subseteq V\left(Q^{n}\right)$ be the set of vertices which are not covered by $M_1$, where we note that $|X| = 2^{n} - 2||\bm{x}||_1 = x_{n+1}$. We form $M_3$ by adding to $M_2$ the edges in direction $\bm{e_{n+1}}$ joining the two copies of each vertex in $X$ in $Q^{n+1}$. That is, 
\[
M_3 := M_2 \cup \{ \big((\bm{v},0), (\bm{v},1) \big) \colon \bm{v} \in X \}.
\]

It is clear that $M_3 \subseteq Q^n$ is a matching with profile $(2\bm{x}, x_{n+1})$ and since $2||\bm{x}||_1 + |X| = 2^{n}$ it follows that $(2\bm{x}, x_{n+1})$ is a  perfect admissible tuple.
\end{proof}

\end{document}
